\documentclass[12pt, a4paper]{article}

\usepackage[margin=2.5cm]{geometry}
\usepackage{amsmath}
\usepackage{amssymb}
\usepackage{mathtools}
\usepackage{setspace}
\usepackage{parskip}
\usepackage{titlesec}
\usepackage{hyperref}
\usepackage{csquotes}
\usepackage{fancyhdr}
\usepackage{microtype}
\usepackage[T1]{fontenc}
\usepackage{lmodern}

\hypersetup{
    colorlinks=true,
    linkcolor=black,
    citecolor=black,
    urlcolor=black
}

\pagestyle{fancy}
\fancyhf{}
\rhead{\small Tribulation of the Precursor}
\lhead{\small GreyFox's Notebook --- Jbaro Archives}
\rfoot{\thepage}

\titleformat{\section}{\large\bfseries}{\thesection}{1em}{}
\titleformat{\subsection}{\normalsize\bfseries}{\thesubsection}{1em}{}

\setlength{\parskip}{0.6em}
\onehalfspacing

\begin{document}

%--- Title block ---
\begin{center}
    {\LARGE\bfseries EXCERPT FROM}\\[0.4em]
    {\LARGE\bfseries TRIBULATION OF THE PRECURSOR}\\[1em]
    {\normalsize\itshape Doing webpages and making apps is not Computer Science,\\
    but a consequence and discipline in the very broad field of Computation}\\[1.5em]
    {\small Summary extract of the collection of private note fragments,\\
    gathered and completed from GreyFox's Notebook}\\[1em]
    {\small Shared for Jbaro Archives, marked by IIX.\\
    January to March 2026\\
    2026 // Open Classification}
\end{center}

\vspace{1.5em}
\hrule
\vspace{1.5em}

%--- Abstract ---
\begin{center}
    {\bfseries ABSTRACT}
\end{center}

Four components of a single inquiry, developed across the period January to March 2026.
\textit{Biocomputers and Uploaded Intelligence} documents the engineering state of the field:
which GPUs have been mapped, which platforms are running biological neural tissue, which algorithms
exist and where they fall short. \textit{The Algorithm of Who We Are} formalizes the identity system
that must remain stable in any emulation attempt for the result to remain the same person, expressed
as a system of differential equations, a fixed-point operator, and an optimal transport condition.
\textit{Computer Science} establishes the epistemic identity and core discipline CS is and why it is
the one that has to solve this. \textit{How Computer Science Borrowed from Neuroscience} traces the
chain of debt from McCulloch in 1943 to the present day.

This excerpt summarizes each component and documents where they bear directly on each other,
including where findings reinforce or contradict one another. No new claims are made. All references
are taken from the source notebooks without modification.

\vspace{1.5em}
\hrule
\vspace{1.5em}

%=============================================================
\section{Biocomputers and Uploaded Intelligence}
%=============================================================

The central constraint of the field is stated at the outset: mapping a connectome and executing it
are two different problems. The connectome alone is not a nervous system: emulating every neuron,
every synapse, every connection. Producing one requires electron microscopy at synaptic resolution,
computational reconstruction, and annotation of cell types. As of 2024, complete connectomes exist
for five organisms. The first was \textit{C.\ elegans}, a nematode with 302 neurons, mapped in 1986
by White, Southgate, Thomson, and Brenner~\cite{Bio-1}. By 2023, complete connectomes existed for
four organisms: \textit{C.\ elegans}, the tadpole larva of \textit{Ciona intestinalis}, the marine
annelid \textit{Platynereis dumerilii}, and the \textit{Drosophila} larva; the last comprising 3,016
neurons and 548,000 synapses, published by Winding et al.\ in \textit{Science}~\cite{Bio-2}. In
2024, the connectome of an entire adult \textit{Drosophila} brain was published by Dorkenwald et al.:
approximately 140,000 neurons and tens of millions of synapses, annotated by cell type and predicted
neurotransmitter identity~\cite{Bio-3}. This is the largest complete adult brain connectome produced
to date.

On the mammalian side, a dataset released in 2025 by the MICrONS Consortium spans approximately one
cubic millimeter of mouse visual cortex, integrating functional imaging of tens of thousands of
neurons with electron microscopy reconstruction containing hundreds of thousands of cells and
approximately half a billion synapses~\cite{Bio-4}. It is not a complete brain. It is a grain of
sand. But it is a grain of sand at synaptic resolution. The scale gap is not a rounding error. It is
the central engineering problem of the field.

Several platforms are currently attempting the execution side of this problem. FinalSpark's
Neuroplatform, operational as of 2024, connects living human cortical organoids to multi-electrode
arrays that allow both electrical stimulation and recording~\cite{Bio-5}. Organoids are clusters of
lab-grown neurons derived from human stem cells that self-organize into layered structures resembling
cortical tissue. Cortical Labs demonstrated in 2022 that cortical neurons cultured on a
multi-electrode array could learn to play a simplified version of Pong using electrophysiological
feedback, without being pre-programmed with the game's rules~\cite{Bio-6}. The neurons adapted their
firing patterns in response to stimulation encoding the game state. The system learned from
experience in a measurable behavioural sense. These platforms do not emulate a whole brain. They
demonstrate that biological neural tissue can be kept alive outside the skull, interfaced with digital
hardware, and made to produce computationally meaningful outputs. That is the engineering foundation
on which any future emulation effort would have to build.

The Sandberg--Bostrom roadmap from 2008, produced through a workshop of invited experts at the Future
of Humanity Institute, Oxford University, identified three core technical requirements for whole brain
emulation: scanning the brain at sufficient resolution to capture the information that determines its
function; translation (converting a structural scan into a computational model, determining which
features are functionally relevant); and simulation (running the model on hardware at a speed and
scale sufficient to produce real-time behaviour)~\cite{Bio-7}. The LOBE research literature also
identifies embodiment as a formal requirement: providing the emulated system with sensory input and
motor output so it can interact with an environment.

The 2025 State of Brain Emulation Report draws the critical formal distinction between simulation and
emulation: simulation matches outputs, emulation reproduces the causal machinery~\cite{Bio-8}. No
current system has demonstrated emulation by that definition for any organism above \textit{C.\
elegans}. The Blue Brain Project ran from 2005 to December 2024 under Henry Markram at EPFL, with the
objective of building a biologically detailed computational model of the neocortical column of a
rat~\cite{Bio-7}. Nineteen years on supercomputing resources. The result was electrophysiological
simulation. Not emulation by Zanichelli's definition.

The distinction between simulation and emulation is not pedantic. It is the precise boundary between
the two questions these materials are asking. The Sandberg--Bostrom roadmap, written in 2008,
anticipated this distinction and stated it clearly. Seventeen years later, no system has crossed from
simulation to emulation for any organism above 302 neurons.

The candidate algorithms for neural simulation exist in published form. The Hodgkin--Huxley model,
published in 1952 and awarded the Nobel Prize in Physiology or Medicine in 1963, is the foundational
algorithm for simulating how individual neurons generate and propagate electrical signals~\cite{Bio-9}.
It expresses the firing mechanism of the action potential as a system of nonlinear differential
equations governing sodium and potassium conductances across the cell membrane. The
Integrate-and-Fire model (Lapicque 1907) is a useful simplification that captures the
threshold-and-fire behaviour without modelling the full ionic dynamics~\cite{Bio-10}. The
Sandberg--Bostrom roadmap describes a hierarchy of simulation levels, from simple rate models up
through full electrophysiological models with protein concentrations and gene expression: higher
integrity requires exponentially more data and compute~\cite{Bio-7}. A single modern GPU can simulate
approximately 0.5 to 1 million neurons, with memory rather than compute as the primary
bottleneck~\cite{Bio-8}. The primary limiting factor is identified as data, not hardware or
algorithms.

\begin{displayquote}
\itshape
``Brain emulation needs to take chemistry more into account than commonly occurs in current
computational models. Chemical processes inside neurons have computational power on their own and
occur on a vast range of timescales, from sub-millisecond to weeks.''\\
\upshape --- Sandberg and Bostrom, \textit{Whole Brain Emulation: A Roadmap}, 2008~\cite{Bio-7}
\end{displayquote}

%=============================================================
\section{The Algorithm of Who We Are}
%=============================================================

This component begins where Biocomputers stops. Connectome data and simulation algorithms address the
structural problem. They do not address the integrity problem: what would need to be preserved for the
result to still be the same person? The component formalises identity as a five-variable system:

\begin{equation}
\mathrm{Identity}(p,\,t) \;:=\; \{\, A(p,t),\; R(p),\; C(p),\; S(p,t),\; M(p,t) \,\}
\end{equation}

where $A$ is attachment state, $R$ is recursive self-reference, $C$ is consciousness, $S$ is the
social network, and $M$ is the autobiographical self. Each has a formal representation grounded in
published research. Each raises a specific open question for the upload scenario.

\subsection*{Attachment (A)}

$A$ is the attachment state. Bowlby's 1969 monograph \textit{Attachment} establishes that the
attachment system is not a learned response to secondary drive but a primary motivational system with
its own neural architecture~\cite{Alg-1}. The neurochemical substrate was documented by Insel and
Young in 2001: oxytocin and vasopressin are the primary mediators of social attachment, operating
through reward circuitry in the nucleus accumbens and ventral tegmental area~\cite{Alg-2}. Fisher,
Aron, and Brown confirmed in 2005 that early-stage romantic love produces a measurable, reproducible
neural signature in the caudate nucleus and VTA, consistent across subjects and distinct from sexual
arousal~\cite{Alg-3}. The governing equation for $A$:

\begin{equation}
\frac{dA}{dt} \;=\; \lambda_{\mathrm{form}} \cdot B(p,t) \cdot [\mathrm{OXT}](p,t)
               \;-\; \lambda_{\mathrm{decay}} \cdot A(t)
               \;-\; \delta(p,t) \cdot A(t)
\end{equation}

where $B(p,t)$ is the dyadic bond function, $[\mathrm{OXT}]$ is oxytocin concentration, and $\delta$
is an external disruption term. When $B(p,t) \to 0$ (the death of an attached person), the equation
reduces to:

\begin{equation}
\frac{dA}{dt} \;=\; -(\lambda_{\mathrm{decay}} + \delta)\,A(t)
\quad\Longrightarrow\quad
A(t) \;=\; A_0\,e^{-(\lambda+\delta)\,t}
\end{equation}

Grief as exponential decay with time constant $\tau = 1/(\lambda+\delta)$. The second-order extension
captures the dynamics of building and fading bonds over time:

\begin{equation}
\frac{d^2A}{dt^2} \;=\; \lambda_{\mathrm{form}}\!\left(\beta_1\,\frac{dB}{dt}
    + \beta_2\,\frac{d[\mathrm{OXT}]}{dt}\right)
    \;-\; \lambda_{\mathrm{decay}}\,\frac{dA}{dt}
\end{equation}

\subsection*{Recursive self-reference (R)}

$R$ is recursive self-reference, formalised as a fixed-point problem on the Hilbert space of
cognitive states. Hofstadter's argument in \textit{G\"{o}del, Escher, Bach} (1979) and \textit{I Am
a Strange Loop} (2007) is that the self is not a thing inside the brain but a process the brain runs:
the process of modelling itself as an object in its own cognitive field~\cite{Alg-4,Alg-5}. The
formal statement:

\begin{equation}
R^*(p) \;=\; F[R^*(p)]
\end{equation}

where $F: L^2(\Omega) \to L^2(\Omega)$ and $\|F(R_1) - F(R_2)\| \leq k\,\|R_1 - R_2\|$,\; $k < 1$.
By the Banach fixed-point theorem, a unique fixed point $R^*$ exists and the iteration
$R_{n+1} = F(R_n)$ converges to it from any starting point. The candidate form for the cognitive
update operator is:

\begin{equation}
F[R](x) \;=\; \int_0^t K(x,x')\,R(x')\,dx' \;+\; \eta(x,\,A(p,t),\,M(p,t))
\end{equation}

where $K(x,x')$ is an attention kernel over the cognitive state space, and $\eta$ is external input
from the attachment and autobiographical memory components. $F$ is a contraction if the operator norm
of $K$ is less than one. Open questions from the notebook: Is $R^*$ unique for each person? Does the
Lipschitz constant $k$ vary with trauma? If the substrate changes at iteration $n$, does the sequence
still converge to the same $R^*$?

\subsection*{Consciousness (C)}

$C$ is consciousness, formalised via Tononi's Integrated Information Theory~\cite{Alg-8}:

\begin{equation}
\Phi(X) \;:=\; \min_{M_k \in P(X)}\; \mathrm{EI}(X \mid M_k)
           \;-\; H(X_t \mid X_{t-1})
           \;-\; H(X_{t-1} \mid X_{M_k})
\end{equation}

$\Phi = 0$ means the system is fully decomposable with no integrated information. $\Phi > 0$ means
irreducible causal structure. Whether any non-biological substrate running a brain map instantiates
the same causal relationships is an open current research question. Baars' Global Workspace theory
(1988) provides a complementary architectural constraint: consciousness requires information to be
broadcast across a global workspace simultaneously available to a wide range of cognitive
processes~\cite{Alg-9}. Chalmers' hard problem (1995) is the deepest obstruction: even after all
functional problems are solved, no test currently exists that can be applied to a candidate upload and
return a reliable answer about whether it is conscious~\cite{Alg-7}.

\subsection*{Social network (S)}

$S$ is the social network, modelled as a weighted graph with Laplacian dynamics:

\begin{equation}
i\hbar\,\frac{\partial w_{ij}}{\partial t}
    \;=\; \gamma \cdot \mathrm{interaction}(i,j,t) \;-\; \mu \cdot w_{ij}(t)
\end{equation}

The graph Laplacian $L = D - W$ yields the Fiedler value $\lambda_2(L)$, algebraic connectivity,
measuring social cohesion. Piolatto et al.\ 2022 confirmed that social isolation worsens cognitive
outcomes through combined mechanisms associated with measurable cognitive decline~\cite{Alg-10}.
Post-upload, old edges freeze or decay: $dw_{ij}/dt = -\mu\,w_{ij}(t)$. The network collapses.

\subsection*{Autobiographical self (M)}

$M$ is the autobiographical self, formalised as a convolution integral over episodic memory, body
state, and experience:

\begin{equation}
M(p,t) \;=\; \int_0^t K(t,\tau)\,\bigl[\,\mathrm{mem}(p,\tau)
           + \mathrm{body}(p,\tau) + \mathrm{exp}(p,\tau)\,\bigr]\,d\tau
\end{equation}

with exponential decay kernel $K(t,\tau) = e^{-\lambda(t-\tau)}$, following the Markov property. The
notebook raises the alternative of a power-law kernel $K \sim (t-\tau)^{-\alpha}$ for long-range
memory. Which kernel is cognitively faithful is an open question with direct computational
implications.

\subsection*{Identity as trajectory}

The identity-as-trajectory formulation is stated explicitly in the source:

\begin{equation}
p(t+dt) \;=\; p(t) \;+\; \frac{dp}{dt}\,dt \;+\; \mathcal{O}(dt^2)
\end{equation}

where $dp/dt = f(A, R, C, S, M, \mathrm{env}(t))$. Identity is a trajectory through state space,
not a point. Uploading copies a snapshot at $t_{\mathrm{upload}}$, not a trajectory. The formal
integrity condition:

\begin{equation}
\mathrm{Integrity}(p_{\mathrm{upload}}) \;:=\; \mathrm{True}
\;\;\text{iff}\;\;
W_2(\mu_{\mathrm{upload}},\, \mu_{\mathrm{orig}}) \;\leq\; \varepsilon
\end{equation}

where $W_2$ is the Wasserstein-2 distance, the minimum cost of transporting one probability
distribution over the identity state space into another. The upload problem is an optimal transport
problem. $W_2 = 0$ requires perfect transport. By the continuity of $W_2$ and the irreversibility of
substrate change, $L > 0$ almost surely. The total loss:

\begin{equation}
L_{\mathrm{total}} \;=\;
    w_A L_A + w_R L_R + w_C L_C + w_S L_S + w_M L_M
\end{equation}

where $L_x = \|x(p_{\mathrm{upload}},t) - x(p_{\mathrm{orig}},t)\|_2$ and the weights $w_x$ are
unknown.

\begin{displayquote}
\itshape
``The equations in these pages are not solutions. They are the shape of the problem, written down so
someone can see it clearly. The connectome gives us the wiring. It does not give us the loop, the
bond, or the light inside.''\\
\upshape --- Algorithm of Who We Are, conclusion
\end{displayquote}

%=============================================================
\section{Computer Science}
%=============================================================

This component establishes the epistemological ground for the entire inquiry: what kind of discipline
CS is, and why the upload problem belongs to it. The argument is that CS occupies a historically
novel position among formal disciplines---simultaneously mathematics, engineering, and empirical
science---as the direct consequence of one structural property that no other formal discipline
possesses at scale.

CS began as formal mathematics. In 1936, Alan Turing published \textit{On Computable Numbers} and
defined a machine that existed only on paper, asking what it could in principle compute~\cite{CS-1}.
Alongside Turing, Alonzo Church developed lambda calculus~\cite{CS-2} and John von Neumann
formalised stored-program architecture~\cite{CS-3}. These were mathematicians drawing boundaries
around a new formal object: computation itself. The name Computer Science was an institutional
accident, applied when universities needed to house a department that fit neither Mathematics nor
Engineering exactly.

The property that changed everything: the Turing machine, a purely theoretical construct, has a
direct physical realisation in the transistor, the chip, the processor. When CS produces a formal
system, that system can be executed. The abstraction does not merely describe a process. It becomes
the process. No other formal discipline has this execution property at scale. The chain runs: logic to
circuit to code to model to body to action.

When CS borrowed the McCulloch--Pitts model and built the Perceptron from it~\cite{CS-4,CS-5}, it
inherited the empirical assumptions of that model. Neural networks are not proven correct through
formal proof. They are validated through testing against data from the real world. When the network
gets a body mounted in a robot with cameras, force sensors, and actuators, the system must be tested
empirically, against the resistance of the real world~\cite{CS-6}. The world pushes back. That is the
defining feature of empirical science.

CS therefore occupies all three categories of rigorous knowledge simultaneously. It is mathematics
when it produces formal proofs about algorithms and computability. It is engineering when it builds
systems that must function under real-world physical constraints. It is empirical science when it
trains models on data and validates them against observable reality. That triple nature is the source
of its extraordinary reach: and it is what makes the upload problem a CS problem. Only CS can attempt
to run a brain map. Only CS has the execution property that would make the attempt more than
theoretical.

\begin{displayquote}
\itshape
``Where would memories go? It becomes a philosophical problem.''\\
\upshape --- GreyFox's Electrical Engineering Professor
\end{displayquote}

%=============================================================
\section{How Computer Science Borrowed from Neuroscience}
%=============================================================

This component traces the historical chain by which CS arrived at the upload frontier. The chain
begins not with an engineer but with two neurophysiologists. In 1943, Warren McCulloch and Walter
Pitts published a paper attempting to describe, in formal logical terms, how a biological neuron
produces an output from its inputs~\cite{Bor-1}. Their answer: the neuron is, in functional terms, a
threshold gate. This was a description of neuroscience, not a proposal for engineering. The paper was
published in the \textit{Bulletin of Mathematical Biophysics}. It was read by engineers.

Frank Rosenblatt at the Cornell Aeronautical Laboratory was one of them. In 1958 he took the
McCulloch--Pitts model and built physical hardware from it~\cite{Bor-2}. The Perceptron received
multiple weighted input signals, summed them, and fired an output when the sum exceeded a threshold.
The design was not derived from mathematical principles of computation. It was lifted directly from
the neurophysiological description of nerve cells. The vocabulary of training, borrowed from animal
learning research, entered CS. The vocabulary of neuroscience followed the hardware into the field
and stayed there.

Minsky and Papert (1969) showed a single-layer perceptron could not solve problems that were not
linearly separable, including the simple XOR operation~\cite{Bor-3}. Funding for connectionist
research largely collapsed through the 1970s. In the 1980s, Rumelhart, Hinton, and Williams
formalised backpropagation: error signals propagated backward through multiple layers, adjusting
weights at each level proportionally~\cite{Bor-4}. The borrowed unit could now be stacked. By the
2010s, deep networks trained on large datasets and run on GPU hardware were outperforming every prior
approach in image recognition, speech recognition, and machine translation~\cite{Bor-5}. Each layer
remained, structurally, a population of threshold units derived from the McCulloch--Pitts neuron.

The Human Connectome Project, NIH-funded since 2009, has been working to chart the structural and
functional connectivity of the human brain using diffusion MRI and functional
neuroimaging~\cite{Bor-6}. FinalSpark's Neuroplatform connects living human cortical organoids to
electrodes for both recording and stimulation~\cite{Bor-8}. DishBrain neurons learned Pong from
feedback without pre-programmed rules~\cite{Bor-9}. These programs close the loop: the direction is
now bidirectional.

Large language models trained on human-generated text were trained on essentially the entire recorded
output of human language and thought. The architecture is biological in origin. The training signal
is human in origin. The result is a system whose outputs were shaped at every stage by something that
came from outside computer science.

\begin{displayquote}
\itshape
``Maybe it's almost as if we are trying to mirror ourselves and something is reflecting back?''\\
\upshape --- GreyFox
\end{displayquote}

%=============================================================
\section{Where the Components Connect}
%=============================================================

These are not four independent arguments. They are one argument approached from four angles, and
reading them together surfaces connections that none of them resolve alone. Several of these
connections constitute direct conflicts or complications that change the meaning of each component
when read in isolation.

The simulation--emulation distinction from Biocomputers~\cite{Bio-8} cuts directly against the $A$
integrity condition from the Algorithm component. The integrity condition is stated as a
Wasserstein-2 distance over identity distributions: $W_2(\mu_{\mathrm{upload}},\mu_{\mathrm{orig}})
\leq \varepsilon$. This metric measures the cost of transporting one probability distribution over
behavioural and cognitive states into another. It is the right metric for identity-as-distribution.
It is not the right metric for identity-as-causal-process. A system can match output distributions
while failing the emulation condition.

The Algorithm component's loss function treats $L_A$ and $L_S$ as independent terms with unknown
weights. The neurobiology cited within the same component does not support this independence. Insel
and Young~\cite{Alg-2} document that attachment state $A$ is mediated by oxytocin and vasopressin
operating through the dopaminergic reward system. Piolatto et al.~\cite{Alg-10} document that social
connectivity $S$ protects cognition through neuroprotective pathways overlapping with the
neurobiological effects of exercise. Both $A$ and $S$ are mediated by the same neurochemical
substrate. Post-upload, the social network $S$ decays as old edge weights follow $dw_{ij}/dt =
-\mu\,w_{ij}$. As $S$ decays, the biological drivers of $A$ lose their input. $L_A$ and $L_S$ are
coupled, not independent. The loss function as stated underestimates the damage that social graph
collapse produces.

The CS component establishes that the only verification method available for any built system is
empirical testing~\cite{Bor-6}. There is a structural incompatibility between what needs to be
verified and the verification framework CS makes available: Chalmers' hard problem~\cite{Alg-7} is
not a philosophical curiosity in this context. It is a direct obstruction to the engineering
verification of the upload.

There is one further connection that reaches across all four components. The Borrowing component ends
with the observation that large language models were trained on the recorded outputs of human minds.
The CS component establishes that CS converts formal objects into systems instantiated in hardware.
The Algorithm component asks what must be preserved for the upload to remain the same person. The
Biocomputers component documents that we are already running human cortical tissue in computational
loops outside the skull. Taken together: CS has been progressively importing more of the brain into
its architecture, then the outputs of human minds, then its neurobiological tissue. The upload
question is not a future scenario. It is the logical conclusion of a trajectory that has been
underway since 1943.

%=============================================================
\section{The Borrowing Trajectory and What It Implies for Uploaded Intelligence}
%=============================================================

The borrowing pattern documented in \textit{How Computer Science Borrowed from Neuroscience} is
toward the upload frontier without explicitly pursuing it. Each step borrowed more of the brain. The
trajectory clarifies where the field actually stands with respect to uploaded intelligence.

The McCulloch--Pitts neuron (1943) was the first loan: the firing behaviour of a single neuron,
abstracted as a threshold gate~\cite{Bor-1}. Nothing about memory, identity, or consciousness was
claimed or attempted. Rosenblatt's Perceptron (1958) borrowed the learning mechanism through which
neural tissue adapts~\cite{Bor-2}. Backpropagation and the multilayer network (1986) borrowed the
cortex's hierarchical architecture~\cite{Bor-4}. Deep learning scaled this to dimensions that have
no biological analogue, running on datasets and compute resources orders of magnitude larger than
anything biological~\cite{Bor-5}. Large language models trained on human-generated text borrowed
something harder to characterise: the outputs of human cognition at scale. Whether the process of
generating that output was captured is precisely the question the Algorithm component asks.

Connectome mapping and biological computing platforms are attempting the next step: borrowing the
wiring. The Human Connectome Project~\cite{Bor-6} maps structural and functional connectivity. The
MICrONS dataset~\cite{Bio-4} reconstructs synaptic-resolution cortical structure. FinalSpark's
Neuroplatform~\cite{Bio-5} and Cortical Labs' DishBrain~\cite{Bio-6,Bor-9} run biological neurons
in computational loops. What is being borrowed at this stage: the actual neural architecture,
connectivity, and tissue of biological brains---not a model of their outputs, but the thing itself.

The progression is: behaviour of one cell, then its learning mechanism, then its hierarchical
architecture, then the outputs of human minds, then the wiring and tissue. The upload would be the
limit of this sequence: not a model of the brain, not a copy of its architecture, not a simulation of
its behaviour, but the brain's causal machinery reproduced with sufficient fidelity that the
Wasserstein-2 distance between the original and the instantiation is within a tolerable
$\varepsilon$. That limit has not been reached. The sequence is ongoing. The direction has not
changed.

%=============================================================
\section{What This Amounts To}
%=============================================================

The question of uploaded intelligence is not a philosophical thought experiment. It is a formal
engineering problem with documented components, published roadmaps, working platforms, and a
mathematically precise integrity condition.

The hardware floor is known and expanding. Complete connectomes exist for five organisms up to
140,000 neurons~\cite{Bio-3}. Mammalian cortical reconstruction exists at synaptic resolution for a
cubic millimetre of mouse visual cortex~\cite{Bio-4}. Living human neural tissue is being interfaced
with digital hardware in real time~\cite{Bio-5,Bio-6}. The simulation algorithms are published and in
active use~\cite{Bio-9,Bio-10}. A single modern GPU can simulate 0.5 to 1 million neurons. The
primary limiting factor is physiological data, not hardware or algorithms~\cite{Bio-8}.

The integrity problem is formally stated. Identity is a five-dimensional trajectory through state
space: attachment $A$, recursive self-reference $R$, consciousness $C$, social network $S$, and
autobiographical self $M$. Three specific blocking problems remain unresolved. The chemistry problem:
no current simulation framework models neurochemical dynamics---specifically the $[\mathrm{OXT}]$
variable in the Algorithm's attachment equation---a data and modelling gap, not a hardware
limitation~\cite{Bio-7}. The consciousness verification problem: $\Phi$ cannot be measured
behaviourally, and the only verification method CS makes available is empirical behavioural
testing~\cite{CS-6}. The trajectory problem: identity is a flow, not a snapshot. An upload captures
a point on a trajectory. Whether the trajectory can restart on a new substrate and converge to the
same attractor depends on whether $F$ is a transferable structural property or an emergent property of
feedback dynamics---a question the DishBrain result~\cite{Bio-6,Bor-9} bears on but does not
resolve.

No current system satisfies the integrity condition. Not for any organism, and certainly not for a
human. The Wasserstein distance between any uploaded mind and its original is greater than zero, and
the tolerable $\varepsilon$ is not known, and neither is who would decide it, or whether the person
on the other side of the upload would be the one who agreed to it. These are not philosophical
complications added after the fact. They follow directly from the formal structure of the integrity
condition as stated in the source material. These are what the mathematics says when you read it
carefully.

\begin{displayquote}
\itshape
``The equations in these pages are not solutions. They are the shape of the problem, written down so
someone can see it clearly. The connectome gives us the loop, the bond, or the light inside.''
\end{displayquote}

\vspace{0.5em}
{\small\itshape Source notebooks: Biocomputers and Uploaded Intelligence; The Algorithm of Who We
Are; Computer Science; How Computer Science Borrowed from Neuroscience. All from GreyFox's private
notebook, shared for Jbaro Archives, January to March 2026.}

%=============================================================
\newpage
\section*{References}
\renewcommand{\section}[2]{}% suppress auto-heading from thebibliography
%=============================================================

{\noindent\bfseries Biocomputers and Uploaded Intelligence}\par\smallskip
\begin{thebibliography}{Bio-10}

\bibitem[Bio-1]{Bio-1} J. G. White, E. Southgate, J. N. Thomson, and S. Brenner, ``The Structure
    of the Nervous System of the Nematode \textit{C. elegans},'' \textit{Philosophical Transactions
    of the Royal Society B}, vol. 314, pp. 1--340, 1986.

\bibitem[Bio-2]{Bio-2} M. Winding et al., ``The Connectome of an Insect Brain,'' \textit{Science},
    vol. 379, no. 6636, 2023.

\bibitem[Bio-3]{Bio-3} S. Dorkenwald et al., ``Neuronal Wiring Diagram of an Adult Brain,''
    \textit{Nature}, vol. 634, pp. 124--138, 2024.

\bibitem[Bio-4]{Bio-4} MICrONS Consortium et al., ``Functional Connectomics Spanning Multiple Areas
    of Mouse Visual Cortex,'' \textit{Nature}, vol. 640, no. 8058, pp. 435--447, 2025.

\bibitem[Bio-5]{Bio-5} F. D. Jordan et al., ``Open and Remotely Accessible Neuroplatform for
    Research in Wetware Computing,'' \textit{Frontiers in Artificial Intelligence}, vol. 7, 2024.

\bibitem[Bio-6]{Bio-6} B. Kagan et al., ``In vitro neurons learn and exhibit sentence when embodied
    in a simulated game-world,'' \textit{Neuron}, vol. 110, no. 23, pp. 3952--3969, 2022.

\bibitem[Bio-7]{Bio-7} A. Sandberg and N. Bostrom, ``Whole Brain Emulation: A Roadmap,'' Future of
    Humanity Institute, Oxford University, Technical Report \#2008-3, 2008.

\bibitem[Bio-8]{Bio-8} N. Zanichelli et al., ``State of Brain Emulation Report 2025,''
    arXiv:2510.15745, 2025.

\bibitem[Bio-9]{Bio-9} A. L. Hodgkin and A. F. Huxley, ``A Quantitative Description of Membrane
    Current and Its Application to Conduction and Excitation in Nerve,'' \textit{Journal of
    Physiology}, vol. 117, pp. 500--544, 1952.

\bibitem[Bio-10]{Bio-10} L. Lapicque, ``Recherches quantitatives sur l'excitation electrique des
    nerfs,'' \textit{Journal de Physiologie et de Pathologie G\'{e}n\'{e}rale}, vol. 9,
    pp. 620--635, 1907.

\end{thebibliography}

\vspace{1em}
{\noindent\bfseries The Algorithm of Who We Are}\par\smallskip
\begin{thebibliography}{Alg-10}

\bibitem[Alg-1]{Alg-1} J. Bowlby, \textit{Attachment and Loss, Vol. 1: Attachment}. New York:
    Basic Books, 1969.

\bibitem[Alg-2]{Alg-2} T. R. Insel and L. J. Young, ``The Neurobiology of Attachment,''
    \textit{Nature Reviews Neuroscience}, vol. 2, no. 2, pp. 129--136, 2001.

\bibitem[Alg-3]{Alg-3} H. Fisher, A. Aron, and L. L. Brown, ``Romantic Love: An fMRI Study of a
    Neural Mechanism for Mate Choice,'' \textit{Journal of Comparative Neurology}, vol. 493, no. 1,
    pp. 58--62, 2005.

\bibitem[Alg-4]{Alg-4} D. R. Hofstadter, \textit{G\"{o}del, Escher, Bach: An Eternal Golden
    Braid}. New York: Basic Books, 1979.

\bibitem[Alg-5]{Alg-5} D. R. Hofstadter, \textit{I Am a Strange Loop}. New York: Basic Books,
    2007.

\bibitem[Alg-6]{Alg-6} A. Damasio, \textit{Self Comes to Mind: Constructing the Conscious Brain}.
    New York: Pantheon Books, 2010.

\bibitem[Alg-7]{Alg-7} D. J. Chalmers, ``Facing Up to the Problem of Consciousness,''
    \textit{Journal of Consciousness Studies}, vol. 2, no. 3, pp. 200--219, 1995.

\bibitem[Alg-8]{Alg-8} G. Tononi, ``An Information Integration Theory of Consciousness,''
    \textit{BMC Neuroscience}, vol. 5, no. 42, 2004.

\bibitem[Alg-9]{Alg-9} B. J. Baars, \textit{A Cognitive Theory of Consciousness}. Cambridge:
    Cambridge University Press, 1988.

\bibitem[Alg-10]{Alg-10} M. Piolatto et al., ``The Effect of Social Relationships on Cognitive
    Decline in Older Adults,'' \textit{BMC Public Health}, vol. 22, no. 278, 2022.

\end{thebibliography}

\vspace{1em}
{\noindent\bfseries Computer Science}\par\smallskip
\begin{thebibliography}{CS-6}

\bibitem[CS-1]{CS-1} A. M. Turing, ``On Computable Numbers, with an Application to the
    Entscheidungsproblem,'' \textit{Proceedings of the London Mathematical Society}, vol. 42,
    pp. 230--265, 1936.

\bibitem[CS-2]{CS-2} A. Church, ``An Unsolvable Problem of Elementary Number Theory,''
    \textit{American Journal of Mathematics}, vol. 58, no. 2, pp. 345--363, 1936.

\bibitem[CS-3]{CS-3} J. von Neumann, ``First Draft of a Report on the EDVAC,'' Moore School of
    Electrical Engineering, University of Pennsylvania, 1945.

\bibitem[CS-4]{CS-4} W. S. McCulloch and W. Pitts, ``A Logical Calculus of the Ideas Immanent in
    Nervous Activity,'' \textit{Bulletin of Mathematical Biophysics}, vol. 5, pp. 115--133, 1943.

\bibitem[CS-5]{CS-5} F. Rosenblatt, ``The Perceptron: A Probabilistic Model for Information
    Storage and Organization in the Brain,'' \textit{Psychological Review}, vol. 65, no. 6,
    pp. 386--408, 1958.

\bibitem[CS-6]{CS-6} K. R. Popper, \textit{The Logic of Scientific Discovery}. London:
    Hutchinson, 1959.

\end{thebibliography}

\vspace{1em}
{\noindent\bfseries How Computer Science Borrowed from Neuroscience}\par\smallskip
\begin{thebibliography}{Bor-10}

\bibitem[Bor-1]{Bor-1} W. S. McCulloch and W. Pitts, ``A Logical Calculus of the Ideas Immanent in
    Nervous Activity,'' \textit{Bulletin of Mathematical Biophysics}, vol. 5, pp. 115--133, 1943.

\bibitem[Bor-2]{Bor-2} F. Rosenblatt, ``The Perceptron: A Probabilistic Model for Information
    Storage and Organization in the Brain,'' \textit{Psychological Review}, vol. 65, no. 6,
    pp. 386--408, 1958.

\bibitem[Bor-3]{Bor-3} M. Minsky and S. Papert, \textit{Perceptrons: An Introduction to
    Computational Geometry}. Cambridge, MA: MIT Press, 1969.

\bibitem[Bor-4]{Bor-4} D. E. Rumelhart, G. E. Hinton, and R. J. Williams, ``Learning
    Representations by Back-Propagating Errors,'' \textit{Nature}, vol. 323, pp. 533--536, 1986.

\bibitem[Bor-5]{Bor-5} Y. LeCun, Y. Bengio, and G. Hinton, ``Deep Learning,'' \textit{Nature},
    vol. 521, pp. 436--444, 2015.

\bibitem[Bor-6]{Bor-6} D. C. Van Essen et al., ``The WU-Minn Human Connectome Project: An
    Overview,'' \textit{NeuroImage}, vol. 80, pp. 62--79, 2013.

\bibitem[Bor-7]{Bor-7} H. Markram et al., ``Reconstruction and Simulation of Neocortical
    Microcircuitry,'' \textit{Cell}, vol. 163, no. 2, pp. 456--492, 2015.

\bibitem[Bor-8]{Bor-8} L. Smirnova et al., ``Organoid Intelligence (OI): The New Frontier in
    Biocomputing and Intelligence-in-a-Dish,'' \textit{Frontiers in Science}, vol. 1, 2023.

\bibitem[Bor-9]{Bor-9} B. Kagan et al., ``In vitro neurons learn and exhibit sentence when embodied
    in a simulated game-world,'' \textit{Neuron}, vol. 110, no. 23, pp. 3952--3969, 2022.

\end{thebibliography}

\end{document}
